\part{The finite maximal-time problem}

\section{The whole-space problem and its first possible singular time}
\label{sec:whole-space-problem}

The three-dimensional Navier--Stokes regularity problem asks whether a smooth
solution can lose its smoothness in finite time.  Fix smooth,
divergence-free, finite-energy initial data \(u_0\), fix \(\nu>0\), and let
\((u,p)\) be the corresponding maximal classical solution on \([0,T_*)\).
We want to prove that \(T_*\) cannot be finite.

Fefferman's formulation of the whole-space Millennium Prize Problem takes
\(u_0\) to be smooth and divergence-free and imposes a decay estimate on every
spatial derivative at every polynomial order.  We write this class as
\begin{equation*}
 \SSigma(\R^3)
 :=\left\{a\in C^\infty(\R^3;\R^3):
 \begin{array}{l}
 \nabla\cdot a=0,\\[-2pt]
 \forall\alpha\in\Nzero^3\ \forall K\in\Nzero\
 \exists C_{\alpha,K}<\infty\ \forall x\in\R^3,\\[-2pt]
 |\partial_x^\alpha a(x)|
 \le C_{\alpha,K}(1+|x|)^{-K}
 \end{array}\right\}.
\end{equation*}
Here \(H^m_\sigma(\R^3):=\{v\in H^m(\R^3;\R^3):\nabla\cdot v=0\}\).  In
particular, \(\SSigma(\R^3)\subset H^m_\sigma(\R^3)\) for every finite \(m\).
The problem asks for velocity and pressure fields satisfying
\begin{equation}
 \partial_tu+(u\cdot\nabla)u+\nabla p=\nu\Delta u,
 \qquad
 \nabla\cdot u=0,
 \qquad
 u(0)=u_0,
 \label{eq:NS}
\end{equation}
with
\[
 u,p\in C^\infty(\R^3\times[0,\infty))
\]
and a kinetic-energy bound uniform in time\bibref{Fefferman}.  Smoothness here
includes the endpoint of every finite interval.  It is stronger than being
smooth on \(\R^3\times(0,T)\) for every \(T<T_*\), because a derivative can
remain finite on every shorter interval and still have no finite trace at
\(T_*\).

The velocity determines the normalized pressure.  Taking the divergence of
the momentum equation and using \(\partial_i u_i=0\) gives
\begin{align*}
 -\Delta p
 &=\partial_i\bigl(u_j\partial_j u_i\bigr)
   =\partial_i\partial_j(u_i u_j).
\end{align*}
On \(\R^3\) we select the Riesz-transform normalization
\begin{equation}
 p=R_iR_j(u_i u_j),
 \qquad
 R_i:=\partial_i(-\Delta)^{-1/2}.
 \label{eq:pressure}
\end{equation}
This removes the additive function of time left by \(\nabla p\) in
\eqref{eq:NS}.  It also fixes an identity that will survive every later
spatial and temporal differentiation.  Pressure is never supplied as an
independent field after the velocity has been chosen.

\subsection{Smoothness, regularity, and continuation}

Three related words will be used throughout the proof, and they should not be
asked to do the same job.  A velocity--pressure pair is \emph{smooth through}
a finite time \(T\) when
\begin{equation*}
 u,p\in C^\infty(\R^3\times[0,T]).
\end{equation*}
Equivalently, every mixed spatial and temporal derivative of finite order
exists on the open spacetime region and extends continuously to both temporal
faces.  Smoothness for all time means smoothness through every finite \(T\).

\emph{Regularity} describes the kind and degree of control presently known.
A Sobolev statement such as \(u(t)\in H^m\), a mixed-norm statement such as
\(u\in L^p_tL^q_x\), and a pointwise statement such as
\(u\in C^k\) are different regularity claims.  None may be silently replaced
by another without the embedding, trace, or equation that connects them.

Suppose a normalized pair \((u,p)\) solves \eqref{eq:NS} on
\([0,T)\).  A \emph{continuation through} \(T\) is a normalized solution of
the same zero-force, fixed-viscosity problem on \([0,T+\varepsilon)\), for
some \(\varepsilon>0\), which agrees with \((u,p)\) before \(T\).  The Clay
problem requires a smooth continuation through every finite time.  A finite
quantity is not enough by itself: a required derivative can remain
finite on every strict subinterval and still fail to approach one value at the
endpoint.

\Needspace{18\baselineskip}
\subsection{Local existence and the maximal history}

An \(H^3\) datum generates a unique pressure-normalized strong solution for a
time depending only on \(\nu\) and its \(H^3\) norm.  Every higher integer
Sobolev order already present in the datum persists on that same interval,
and a finite maximal lifespan forces the \(H^3\) norm to escape.  These are
the local facts used at both ends of the proof\bibref{FujitaKato,Kato}; the
proposition states them together, and Appendix~\ref{app:local-theory} gives
the construction in the form required below.

The precise statement is the following.

\begin{proposition}[Maximal history and blow-up alternative]
\label{prop:maximal-history}
Let \(t_0\ge0\) and \(a\in H^3_\sigma(\R^3)\).  There is a time
\(\delta=\delta(\nu,\norm{a}_{H^3})>0\) and a unique pressure-normalized
solution \((w,q)\) of \eqref{eq:NS} in the class
\begin{equation}
 \begin{aligned}
 w&\in C([t_0,t_0+\delta];H^3_\sigma)
      \cap L^2(t_0,t_0+\delta;H^4),\\
 w_t&\in L^2(t_0,t_0+\delta;H^2),
 \qquad
 q=R_iR_j(w_iw_j).
 \end{aligned}
 \label{eq:local-base-class}
\end{equation}
with \(w(t_0)=a\).  If \(a\in H^m_\sigma\) for an integer \(m>3\), then, on
the same interval,
\begin{equation*}
 w\in C([t_0,t_0+\delta];H^m_\sigma)
   \cap L^2(t_0,t_0+\delta;H^{m+1}),
 \qquad
 w_t\in L^2(t_0,t_0+\delta;H^{m-1}).
\end{equation*}
Compatible local solutions in these strong classes join by
uniqueness into one maximal solution on \([t_0,T_{\max})\).  If
\(T_{\max}<\infty\), then
\begin{equation*}
 \limsup_{t\uparrow T_{\max}}\norm{w(t)}_{H^3}=\infty.
\end{equation*}
\end{proposition}

Suppose instead that
\(\norm{w(t)}_{H^3}\le M\) for all \(t<T_{\max}\).  The local construction
then supplies the same positive lifespan
\(\delta(\nu,M)\) whenever it is restarted from any time \(t<T_{\max}\).
Choose
\(t>T_{\max}-\tfrac12\delta(\nu,M)\).  The restarted solution crosses
\(T_{\max}\), and uniqueness identifies it with \(w\) on the overlap.  This
extends the maximal solution.  The extension uses the actual value of the
same solution at a late time; it does not select a new terminal datum.

\Needspace{11\baselineskip}
For \(u_0\in\SSigma(\R^3)\), Proposition~\ref{prop:maximal-history} produces
a smooth normalized solution on a short interval.  Repeated application and
uniqueness produce one maximal pair \((u,p)\) on \([0,T_*)\).  Every finite
spatial order present at time zero persists on each compact subinterval
\([0,T]\subset[0,T_*)\).  Equation~\eqref{eq:NS} and the pressure identity
then recover all finite temporal derivatives there.  Thus no regularity is
missing before \(T_*\).  The unresolved distinction is between bounds on
every fixed \([0,T]\), whose constants may diverge as \(T\uparrow T_*\), and
one bound that survives on the full candidate interval \((0,T_*)\).
Proposition~\ref{prop:maximal-history} makes that distinction exact: a
whole-terminal \(H^3\) bound supplies a uniform restart time and excludes
finite maximality.

The global assertion is the following.

\begin{theorem}[Whole-space existence and smoothness]
\label{thm:main}
For every \(\nu>0\) and every \(u_0\in\SSigma(\R^3)\), the unforced
three-dimensional incompressible Navier--Stokes equations have a global
pressure-normalized classical solution satisfying
\[
 u,p\in C^\infty(\R^3\times[0,\infty)).
\]
For every \(t\ge0\), this solution obeys
\[
 \frac12\norm{u(t)}_2^2
 +\nu\int_0^t\norm{\nabla u(\tau)}_2^2\,d\tau
 =\frac12\norm{u_0}_2^2,
\]
and in particular
\begin{equation}
 \sup_{t\ge0}\int_{\R^3}|u(x,t)|^2\,dx
 \le \norm{u_0}_2^2<\infty.
 \label{eq:uniform-kinetic-energy}
\end{equation}
\end{theorem}

The energy identity will give the uniform kinetic-energy bound in
Theorem~\ref{thm:main}.  Its harder conclusion begins one derivative higher.
Even if the continuation-grade \(H^3\) norm escapes, the kinetic energy and an
integrated first derivative remain controlled for the entire proposed branch.
Those are the first quantities that do not end when the classical restart norm
does.  We derive them next and ask how much of the original evolution they
still identify.

\section{The first global control: energy and weak continuation}
\label{sec:energy-weak}

Pair the momentum equation with \(u\) and integrate over \(\R^3\).  The
transport, pressure, and viscous terms have three distinct forms, so we record
the three integrations separately:
\begin{align*}
 \int_{\R^3}(u\cdot\nabla)u\cdot u\,dx
 &=\frac12\int_{\R^3}u\cdot\nabla|u|^2\,dx
 =-\frac12\int_{\R^3}(\nabla\cdot u)|u|^2\,dx=0,\\
 \int_{\R^3}\nabla p\cdot u\,dx
 &=-\int_{\R^3}p\,\nabla\cdot u\,dx=0,\\
 \int_{\R^3}\Delta u\cdot u\,dx
 &=-\int_{\R^3}|\nabla u|^2\,dx.
\end{align*}
The decay inherited on compact time intervals justifies the integrations;
the same calculation follows by cutoff and density at finite energy.  Hence
\begin{equation*}
 \frac12\frac d{dt}\norm{u(t)}_2^2
 +\nu\norm{\nabla u(t)}_2^2=0.
\end{equation*}
Integrating from \(0\) to \(t<T_*\) gives the exact energy identity
\begin{equation}
 \frac12\norm{u(t)}_2^2
 +\nu\int_0^t\norm{\nabla u(\tau)}_2^2\,d\tau
 =\frac12\norm{u_0}_2^2.
 \label{eq:energy}
\end{equation}

The two controls in \eqref{eq:energy} have different strengths in time.  The
velocity has a pointwise-in-time \(L^2\) ceiling.  Its first spatial derivative
is controlled only after integration over the interval:
\begin{equation}
 \sup_{0\le t<T_*}\norm{u(t)}_2^2\le\norm{u_0}_2^2,
 \qquad
 \int_0^{T_*}\norm{\nabla u(t)}_2^2\,dt
 \le\frac{\norm{u_0}_2^2}{2\nu}
 \quad(T_*<\infty).
 \label{eq:energy-class-controls}
\end{equation}
A finite integral permits high values on intervals whose lengths shrink fast
enough.  The energy law rules out loss of total kinetic energy control, but it
does not place an instantaneous ceiling on the derivative that controls
classical continuation.

\subsection{What global weak existence does and does not continue}

Leray used the bounds in \eqref{eq:energy-class-controls} to construct a
global weak solution.  For the same datum \(u_0\), there is a Leray--Hopf pair
\((U,P)\) on every finite time interval with
\begin{equation*}
 U\in L^\infty(0,T;L^2_\sigma)
    \cap L^2(0,T;H^1_\sigma),
\end{equation*}
which solves \eqref{eq:NS} in distributions and satisfies the energy
inequality\bibref{Leray}.

Before \(T_*\), this weak solution is the classical solution already fixed
above.  To see the identification directly, fix \(T<T_*\), set \(z=U-u\),
and subtract the two equations.  Using
\[
 (U\cdot\nabla)U-(u\cdot\nabla)u
 =(U\cdot\nabla)z+(z\cdot\nabla)u,
\]
pair the difference equation with \(z\).  The first transport term cancels
because \(\nabla\cdot U=0\), and the pressure difference cancels because
\(\nabla\cdot z=0\).  The remaining nonlinear term satisfies
\[
 \left|\int_{\R^3}(z\cdot\nabla)u\cdot z\,dx\right|
 \le \norm{\nabla u}_\infty\norm{z}_2^2.
\]
The difference energy inequality is
\begin{align*}
 \frac12\norm{z(t)}_2^2
 +\nu\int_0^t\norm{\nabla z(s)}_2^2\,ds
 &\le
 \int_0^t\norm{\nabla u(s)}_\infty\norm{z(s)}_2^2\,ds.
\end{align*}
The classical solution belongs to \(C([0,T];H^3)\).  The embedding
\(H^3(\R^3)\hookrightarrow W^{1,\infty}(\R^3)\), proved in
Appendix~\ref{app:analytic-framework} at
\eqref{eq:appendix-spatial-embedding}, gives
\(\norm{\nabla u}_\infty\in L^1(0,T)\).  Since \(z(0)=0\), Gronwall's
inequality gives \(z=0\) on \([0,T]\)\bibref{Prodi,SerrinIVP}.  As \(T<T_*\)
was arbitrary,
\begin{equation}
 U=u\quad\hbox{almost everywhere on }\R^3\times(0,T_*).
 \label{eq:weak-strong-identification}
\end{equation}

Thus a finite \(T_*\) would not prevent a Leray--Hopf continuation from
existing.  That continuation agrees with the fixed classical solution before
\(T_*\), but uniqueness beyond the endpoint has not been established.  The
loss at \(T_*\) would be the loss of a smooth state that identifies one unique
classical continuation, not the disappearance of every distributional
solution.  The next estimate asks where the first classical derivative can
escape while the energy law remains intact.

\section{The three-dimensional obstruction in the first derivative}
\label{sec:enstrophy}

The kinetic-energy identity does not expose the mechanism that makes the
three-dimensional problem difficult.  The first derivative does.  Define the
vorticity
\[
 \omega:=\nabla\times u.
\]
For a decaying vector field,
\[
 \int_{\R^3}|\nabla u|^2\,dx
 =\int_{\R^3}\bigl(|\nabla\times u|^2+|\nabla\cdot u|^2\bigr)\,dx,
\]
so incompressibility gives
\begin{equation*}
 \norm{\nabla u}_2^2=\norm{\omega}_2^2.
\end{equation*}
The dissipation in the kinetic-energy law is exactly the time-integrated
enstrophy.

Taking the curl of the momentum equation removes the pressure gradient.  The
transport term does not pass through the curl unchanged.  In index notation,
using \(\partial_j u_j=0\),
\[
 \nabla\times\bigl((u\cdot\nabla)u\bigr)
 =(u\cdot\nabla)\omega-(\omega\cdot\nabla)u.
\]
The vorticity consequently satisfies
\begin{equation}
 \partial_t\omega+(u\cdot\nabla)\omega
 =(\omega\cdot\nabla)u+\nu\Delta\omega.
 \label{eq:vorticity}
\end{equation}
Transport carries vorticity through the velocity field, viscosity diffuses
its spatial variation, and \((\omega\cdot\nabla)u\) changes vorticity through
the velocity gradient already present in that field.

Write
\[
 S:=\frac12\bigl(\nabla u+(\nabla u)^{\mathsf T}\bigr),
 \qquad
 A:=\frac12\bigl(\nabla u-(\nabla u)^{\mathsf T}\bigr).
\]
Pairing \eqref{eq:vorticity} with \(\omega\) cancels transport and integrates
the Laplacian by parts:
\[
 \frac12\frac d{dt}\norm{\omega}_2^2
 +\nu\norm{\nabla\omega}_2^2
 =\int_{\R^3}\bigl((\omega\cdot\nabla)u\bigr)\cdot\omega\,dx.
\]
Since \((A\omega)\cdot\omega=0\), this is
\begin{equation*}
 \frac12\frac d{dt}\norm{\omega}_2^2
 +\nu\norm{\nabla\omega}_2^2
 =\int_{\R^3}(S\omega)\cdot\omega\,dx.
\end{equation*}
The right-hand side has no fixed sign.  The strain may shorten vorticity in
one region and lengthen it in another, with both actions constrained by the
same divergence-free velocity.  This fully three-dimensional stretching term
prevents the enstrophy balance from inheriting the monotonic form of the
kinetic-energy balance.

The absence of an explicit pressure term from \eqref{eq:vorticity} does not
turn this into a local scalar equation.  Since \(\nabla\cdot u=0\),
\[
 -\Delta u=\nabla\times\omega,
 \qquad
 \partial_k u_j
 =\partial_k(-\Delta)^{-1}(\nabla\times\omega)_j.
\]
Thus every component of \(\nabla u\) is an order-zero singular-integral
transform of \(\omega\); the required multiplier bounds are
\eqref{eq:appendix-riesz-multiplier-bounds}.  The normalized pressure is recovered
from \(u\) by \eqref{eq:pressure}.  Stretching, transport,
viscous diffusion, pressure response, and incompressibility remain parts of
the same evolving field.  The point of the vorticity equation is narrower:
it locates the sign-indefinite derivative production that the total energy
balance cannot see.

\subsection{The first differentiated estimate}

Set
\[
 Y(t):=\norm{\nabla u(t)}_2^2.
\]
Pair \eqref{eq:NS} with \(-\Delta u\).  The pressure vanishes because
\(\nabla\cdot\Delta u=0\), and integration by parts in the time and viscous
terms gives
\begin{equation}
 \frac12Y'(t)+\nu\norm{\Delta u(t)}_2^2
 =\int_{\R^3}(u\cdot\nabla)u\cdot\Delta u\,dx.
 \label{eq:first-differentiated-energy}
\end{equation}
The transport term now survives.  Hölder's inequality with exponents
\((6,3,2)\), followed by the homogeneous Sobolev and interpolation estimates
proved in \eqref{eq:appendix-homogeneous-Sobolev-interpolation}, gives the
classical enstrophy estimate:
\begin{align*}
 \left|\int_{\R^3}(u\cdot\nabla)u\cdot\Delta u\,dx\right|
 &\le \norm{u}_6\norm{\nabla u}_3\norm{\Delta u}_2\\
 &\le C\norm{\nabla u}_2
          \bigl(\norm{\nabla u}_2^{1/2}
                \norm{\nabla u}_6^{1/2}\bigr)
          \norm{\Delta u}_2\\
 &\le C\norm{\nabla u}_2^{3/2}
          \norm{\Delta u}_2^{3/2}.
\end{align*}
Young's inequality, with the viscosity retained in the absorbing term, yields
\[
 C\norm{\nabla u}_2^{3/2}\norm{\Delta u}_2^{3/2}
 \le \frac{\nu}{2}\norm{\Delta u}_2^2
      +C\nu^{-3}\norm{\nabla u}_2^6.
\]
Substitution in \eqref{eq:first-differentiated-energy} gives
\begin{equation}
 Y'(t)+\nu\norm{\Delta u(t)}_2^2
 \le C\nu^{-3}Y(t)^3.
 \label{eq:enstrophy-obstruction}
\end{equation}

The energy identity controls \(\int_0^T Y(t)\,dt\).  The next estimate asks
for \(\int_0^T Y(t)^3\,dt\).  These are genuinely different demands.  A
nonnegative function can have finite integral while its cube has infinite
integral, and the energy identity contains no additional time information
that excludes that pattern.  This observation does not assert that a
Navier--Stokes solution can realize an arbitrary spike.  It identifies the
exact point at which the basic estimate stops: the available \(L^1_t\) control
of \(Y\) cannot be inserted into the cubic term of
\eqref{eq:enstrophy-obstruction}.

\Needspace{9\baselineskip}
The same differential inequality forces a minimum growth rate on any branch
where enstrophy becomes unbounded.  A nonzero decaying incompressible solution
has \(Y(t)>0\) along such a branch.  If \(Y(t_0)=0\), then \(u(t_0)=0\), and
forward uniqueness forces \(u(t)=0\) for every \(t\ge t_0\), excluding any
later blow-up.  Dropping the nonnegative viscous term from
\eqref{eq:enstrophy-obstruction} and differentiating \(Y^{-2}\) gives
\[
 \frac d{dt}Y^{-2}
 =-2Y^{-3}Y'
 \ge -2C\nu^{-3}.
\]
Assume \(T_*<\infty\) and choose \(s_j\uparrow T_*\) with
\(Y(s_j)\to\infty\).  Integration from \(t\) to \(s_j\) gives
\[
 Y(t)^{-2}
 \le Y(s_j)^{-2}+2C\nu^{-3}(s_j-t).
\]
Letting \(j\to\infty\) yields
\begin{equation}
 T_*<\infty\ \hbox{and}\
 \limsup_{t\uparrow T_*}Y(t)=\infty
 \quad\Longrightarrow\quad
 Y(t)\ge
 \frac{\nu^{3/2}}{\sqrt{2C}}\,(T_*-t)^{-1/2}.
 \label{eq:leray-enstrophy-rate}
\end{equation}
The rate in \eqref{eq:leray-enstrophy-rate} is unbounded and still integrable
at \(T_*\).  It can fit inside the dissipation budget only because the time
available above each higher level contracts.  The first differentiated
estimate has consequently exposed a possible terminal geometry: increasing
derivative intensity supported on shrinking time intervals.  Since \(Y(t)\)
is integrated over all of \(\R^3\), it says nothing yet about where that
intensity sits or what remains when the surrounding spacetime region is
magnified.

\subsection{From shrinking time to shrinking spacetime}
\label{sec:CKN}

\Needspace{10\baselineskip}
That missing spatial information is exactly what the local energy inequality
was built to retain.  Caffarelli, Kohn, and Nirenberg inserted a nonnegative
spacetime cutoff into the Leray energy argument, keeping both the energy
inside a shrinking region and the flux through its boundary.  The Leray pair
fixed in Section~\ref{sec:energy-weak} may be chosen suitable; it retains the
same normalized pressure and still satisfies
\eqref{eq:weak-strong-identification} before \(T_*\).  This pair \((U,P)\)
satisfies, for every nonnegative
\(\phi\in C_c^\infty(\R^3\times[0,T))\) and almost every \(t\in(0,T)\),
\begin{equation}
\begin{aligned}
 \frac12\int_{\R^3}|U(t)|^2\phi(t)\,dx
 &+\nu\int_0^t\int_{\R^3}|\nabla U|^2\phi\,dx\,ds\\
 &\le \frac12\int_{\R^3}|u_0|^2\phi(0)\,dx\\
 &\quad+\frac12\int_0^t\int_{\R^3}
 |U|^2(\partial_s\phi+\nu\Delta\phi)\,dx\,ds\\
 &\quad+\int_0^t\int_{\R^3}
 \left(\frac12|U|^2+P\right)U\cdot\nabla\phi\,dx\,ds.
\end{aligned}
\label{eq:local-energy-inequality}
\end{equation}
Unlike the total energy law, \eqref{eq:local-energy-inequality} retains the
flux through the region selected by \(\phi\).  It is designed to detect a
point around which energy and dissipation keep concentrating as the region
shrinks.

For \(z_0=(x_0,t_0)\), write
\[
 Q_r(z_0):=B_r(x_0)\times(t_0-r^2,t_0).
\]
Under the Navier--Stokes scaling
\[
 U_\lambda(x,t)=\lambda U(\lambda x,\lambda^2t),
 \qquad
 P_\lambda(x,t)=\lambda^2P(\lambda x,\lambda^2t),
\]
a spatial length \(r\) is paired with a time length \(r^2\).  If
\(z_0^\lambda=(x_0/\lambda,t_0/\lambda^2)\), then the change of variables
\(y=\lambda x\), \(s=\lambda^2t\) gives
\[
 \frac{1}{r/\lambda}
 \int_{Q_{r/\lambda}(z_0^\lambda)}|\nabla U_\lambda|^2\,dx\,dt
 =\frac1r\int_{Q_r(z_0)}|\nabla U|^2\,dy\,ds.
\]
Thus the scaled dissipation used below is genuinely dimensionless.  In the
Caffarelli--Kohn--Nirenberg terminology, \(z_0\) is a regular point when \(U\)
is essentially bounded on a spacetime neighbourhood of \(z_0\).  Their
epsilon-regularity theorem gives a constant
\(\varepsilon_{\mathrm{CKN}}(\nu)>0\) such that\bibref{CKN}
\begin{equation}
 \limsup_{r\downarrow0}\frac1r
 \int_{Q_r(z_0)}|\nabla U|^2\,dx\,dt
 <\varepsilon_{\mathrm{CKN}}(\nu)
 \quad\Longrightarrow\quad
 z_0\ \hbox{is a regular point}.
 \label{eq:CKN-epsilon}
\end{equation}
At a singular point the contrapositive of \eqref{eq:CKN-epsilon} gives
\[
 \limsup_{r\downarrow0}\frac1r
 \int_{Q_r(z_0)}|\nabla U|^2\,dx\,dt
 \ge\varepsilon_{\mathrm{CKN}}(\nu).
\]
Thus the scale-invariant dissipation reaches the threshold along arbitrarily
small cylinders.  The limsup statement does not assert the same lower bound
at every small radius.

\Needspace{7\baselineskip}
The covering argument turns this scale-by-scale lower concentration into a
geometric restriction.  If \(\mathcal S(U)\) denotes the set of singular
spacetime points of a suitable weak solution, then
\begin{equation*}
 \mathcal P^1\bigl(\mathcal S(U)\bigr)=0,
\end{equation*}
where \(\mathcal P^1\) is one-dimensional parabolic Hausdorff
measure\bibref{CKN}.  This is far stronger than an almost-everywhere statement.
A possible singular set cannot contain a parabolic curve segment, a positive
time interval, or any set carrying positive one-dimensional parabolic
measure.

Caffarelli--Kohn--Nirenberg have now reduced a possible failure from diffuse
loss of control to scale-invariant concentration near a very small singular
set.  That geometric restriction still leaves an isolated singular point
available, and it still supplies no whole-space state at \(T_*\).  On
\(\R^3\), concentration may also move through space, so regularity in each
fixed neighbourhood does not by itself bound a global Sobolev norm.  The
analysis has reached a precise divide: local concentration describes how failure
would appear, while continuation requires a global quantity strong enough to
propagate the complete spatial state.

\section{What would be strong enough to continue the solution}
\label{sec:high-order-continuation}

Proposition~\ref{prop:maximal-history} already names one such state: a bounded
\(H^3\) norm crosses the proposed endpoint.  Differentiating the equation
shows something stronger.  Once one lower-order transport coefficient is
integrable, every fixed Sobolev order propagates on the same interval.  Let
\(m\ge3\) be an integer and apply \(\partial^\alpha\) to \eqref{eq:NS} for
\(|\alpha|\le m\):
\[
 \partial_t\partial^\alpha u
 +(u\cdot\nabla)\partial^\alpha u
 +[\partial^\alpha,u\cdot\nabla]u
 +\nabla\partial^\alpha p
 =\nu\Delta\partial^\alpha u.
\]
Pair with \(\partial^\alpha u\) and integrate.  The undifferentiated transport
coefficient still cancels,
\[
 \int_{\R^3}(u\cdot\nabla)\partial^\alpha u
             \cdot\partial^\alpha u\,dx=0,
\]
and the pressure cancels because
\(\nabla\cdot\partial^\alpha u=0\).  The commutator contains all the product
terms in which at least one derivative lands on the transport field:
\[
 [\partial^\alpha,u\cdot\nabla]u
 =\sum_{\substack{0<\beta\le\alpha}}
   \binom{\alpha}{\beta}
   \bigl(\partial^\beta u\cdot\nabla\bigr)
   \partial^{\alpha-\beta}u.
\]
The integer-order commutator estimate
\eqref{eq:appendix-commutator} gives
\[
 \sum_{|\alpha|\le m}
 \left|
  \pair{[\partial^\alpha,u\cdot\nabla]u}{\partial^\alpha u}
 \right|
 \le C_m\norm{\nabla u}_\infty\norm{u}_{H^m}^2.
\]
Summing over \(|\alpha|\le m\) produces
\begin{equation}
 \frac12\frac d{dt}\norm{u(t)}_{H^m}^2
 +\nu\norm{\nabla u(t)}_{H^m}^2
 \le
 C_m\norm{\nabla u(t)}_\infty\norm{u(t)}_{H^m}^2.
 \label{eq:high-order-continuation}
\end{equation}

Gronwall's inequality now states the continuation mechanism in one line:
\begin{equation}
 \norm{u(t)}_{H^m}^2
 \le \norm{u_0}_{H^m}^2
 \exp\left(
  2C_m\int_0^t\norm{\nabla u(s)}_\infty\,ds
 \right),
 \qquad 0\le t<T_*.
 \label{eq:high-order-gronwall}
\end{equation}
The factor \(\norm{u_0}_{H^m}\) is finite at every fixed order because
\(u_0\in\SSigma\).  The constant may depend on \(m\); smoothness does not
require one constant uniform over all derivative orders.  What must be common
to every row is the same finite time interval and the same integral
\[
 \int_0^{T_*}\norm{\nabla u(t)}_\infty\,dt.
\]
If this integral is finite, \eqref{eq:high-order-gronwall} bounds every fixed
\(H^m\) norm up to \(T_*\), beginning with \(H^3\).  Proposition
\ref{prop:maximal-history} then continues the solution.  Hence
\begin{equation}
 T_*<\infty
 \quad\Longrightarrow\quad
 \int_0^{T_*}\norm{\nabla u(t)}_\infty\,dt=\infty.
 \label{eq:lipschitz-blowup-criterion}
\end{equation}

A single uniform Sobolev bound above the Lipschitz threshold also closes the
argument.  If \(s>5/2\) and
\[
 \sup_{t<T_*}\norm{u(t)}_{H^s}<\infty,
\]
then \(H^s(\R^3)\hookrightarrow W^{1,\infty}(\R^3)\).  The finite length of
\((0,T_*)\) makes the coefficient in
\eqref{eq:high-order-continuation} integrable, all higher orders propagate,
and Proposition~\ref{prop:maximal-history} restarts the solution.  Thus every
finite maximal time
must force every continuation-grade Sobolev norm to escape.

The calculation has isolated one coefficient that carries the entire spatial
ladder.  To see how far the energy class lies from paying it, first observe
that energy already reaches the spatial exponent that will become critical,
but with the wrong time control.  Sobolev embedding and interpolation
give
\begin{equation*}
 \norm{u(t)}_3
 \le \norm{u(t)}_2^{1/2}\norm{u(t)}_6^{1/2}
 \le C\norm{u(t)}_2^{1/2}\norm{\nabla u(t)}_2^{1/2}.
\end{equation*}
Raising to the fourth power, integrating in time, and using
\eqref{eq:energy} yields
\begin{equation*}
 u\in L^4(0,T_*;L^3(\R^3))
 \qquad (T_*<\infty).
\end{equation*}
Energy already reaches the spatial exponent \(3\), but only after time
integration.  The scale-balanced criteria distinguish exactly how much
stronger the temporal control must be.

\subsection{Scale-balanced continuation and the critical endpoint}
\label{sec:LPS}

Ladyzhenskaya, Prodi, and Serrin found the first scale-balanced way to obtain
that control.  Their mixed space--time conditions trade spatial intensity
against the length of time for which that intensity persists
\bibref{Prodi,SerrinIVP,Ladyzhenskaya,Serrin}.  The first differentiated
estimate determines the trade, so its exponents are derived before the
criterion is stated.

Fix \(q>3\) and choose \(r=2q/(q-2)\), so
\[
 \frac1q+\frac1r+\frac12=1.
\]
\Needspace{10\baselineskip}
Hölder's inequality and the interpolation formula in
\eqref{eq:appendix-homogeneous-Sobolev-interpolation} give
\begin{align*}
 \left|\int_{\R^3}(u\cdot\nabla)u\cdot\Delta u\,dx\right|
 &\le \norm{u}_q\norm{\nabla u}_r\norm{\Delta u}_2,\\
 \norm{\nabla u}_r
 &\le C\norm{\nabla u}_2^{1-3/q}
          \norm{\nabla u}_6^{3/q}\\
 &\le C\norm{\nabla u}_2^{1-3/q}
          \norm{\Delta u}_2^{3/q}.
\end{align*}
Thus
\[
 \left|\int_{\R^3}(u\cdot\nabla)u\cdot\Delta u\,dx\right|
 \le
 C\norm{u}_q
 \norm{\nabla u}_2^{1-3/q}
 \norm{\Delta u}_2^{1+3/q}.
\]
Young's inequality uses the exponent \(2/(1+3/q)\) on the last factor and
its conjugate \(2/(1-3/q)=2q/(q-3)\) on the remaining product.  It follows
that
\begin{equation}
 \frac d{dt}\norm{\nabla u}_2^2
 +\nu\norm{\Delta u}_2^2
 \le
 C_{\nu,q}\norm{u}_q^{\,2q/(q-3)}
 \norm{\nabla u}_2^2.
 \label{eq:LPS-differential}
\end{equation}
The exponent is not chosen for appearance.  It is exactly the time exponent
that makes the coefficient in \eqref{eq:LPS-differential} integrable on the
critical line.

\begin{proposition}[Ladyzhenskaya--Prodi--Serrin continuation criterion]
\label{prop:LPS}
Let \((u,p)\) be the maximal solution supplied by
Proposition~\ref{prop:maximal-history}, and let \(T\le T_*\).
If
\begin{equation*}
 u\in L^p(0,T;L^q(\R^3)),
 \qquad
 \frac2p+\frac3q\le1,
 \qquad
 q>3,
\end{equation*}
then \(u\) remains regular through \(T\).  The endpoint
\((p,q)=(2,\infty)\) is included.
\end{proposition}

The proposition is the classical Ladyzhenskaya--Prodi--Serrin continuation
theorem\bibref{Ladyzhenskaya,Prodi,Serrin}.  The calculation above explains
its critical exponent in the present notation: \(2/p+3/q\le1\) gives
\(p\ge2q/(q-3)\), so on a finite interval the coefficient
\(\norm{u}_q^{2q/(q-3)}\) in \eqref{eq:LPS-differential} is integrable.
Gronwall bounds enstrophy and the integrated \(H^2\) dissipation; the cited
theorem supplies continuation and weak--strong uniqueness through the
terminal time.

The continuation criterion also determines when the lower rate
\eqref{eq:leray-enstrophy-rate} applies.  If
\(\sup_{t<T_*}\norm{\nabla u(t)}_2<\infty\), Sobolev embedding gives
\[
 \sup_{t<T_*}\norm{u(t)}_6<\infty.
\]
Since \(T_*<\infty\), this places \(u\) in \(L^4(0,T_*;L^6)\), and
\[
 \frac24+\frac36=1.
\]
Proposition~\ref{prop:LPS} would continue the solution.  A finite maximal time
must consequently have unbounded enstrophy, so the Leray lower rate applies.

\subsection{Why the endpoint is critical}

The Navier--Stokes scaling is
\begin{equation}
 u_\lambda(x,t)=\lambda u(\lambda x,\lambda^2t),
 \qquad
 p_\lambda(x,t)=\lambda^2p(\lambda x,\lambda^2t).
 \label{eq:NS-scaling}
\end{equation}
At corresponding times and intervals,
\[
 \norm{u_\lambda}_{L^p_tL^q_x}
 =\lambda^{\,1-2/p-3/q}\norm{u}_{L^p_tL^q_x}.
\]
The equality \(2/p+3/q=1\) is precisely the scale-invariant line.  The
spatial exponent \(q\) measures the intensity of concentration, while the
temporal exponent \(p\) measures how long that intensity persists.  The
\Needspace{5\baselineskip}
Leray energy controls
\[
 u\in L^\infty_tL^2_x\cap L^2_tL^6_x,
\]
and both pairs have \(2/p+3/q=3/2\), above the continuation line.  Scaling
expresses the same gap seen in the enstrophy calculation: the global energy
class does not forbid activity that becomes stronger as its scale contracts.

As \(q\downarrow3\), the exponent \(2q/(q-3)\) in
\eqref{eq:LPS-differential} diverges.  The direct Gronwall calculation reaches
every \(q>3\) and stops at
\[
 u\in L^\infty(0,T;L^3(\R^3)).
\]
The norm is exactly invariant under \eqref{eq:NS-scaling}:
\[
 \norm{u_\lambda(t)}_3
 =\norm{u(\lambda^2t)}_3.
\]
At the same time,
\[
 \norm{u_\lambda(t)}_2^2
 =\lambda^{-1}\norm{u(\lambda^2t)}_2^2.
\]
Concentration to a shorter spatial scale preserves the \(L^3\) size while it
reduces the amount seen by the \(L^2\) energy.  This is why the endpoint cannot
be obtained by interpolating the energy law alone.

\subsection{Escauriaza--Seregin--\v{S}ver\'ak}

Escauriaza, Seregin, and \v{S}ver\'ak proved that boundedness at the missing
endpoint is still enough for regularity\bibref{ESS}.  In the continuation form
used here, their theorem states that a suitable Leray--Hopf solution on
\(\R^3\times(0,T)\) satisfying
\[
 U\in L^\infty(0,T;L^3(\R^3))
\]
is regular through \(T\).  The proof cannot follow
\eqref{eq:LPS-differential}, whose exponent has become infinite.  It rescales
around a presumed concentration, passes to a nontrivial ancient limiting
solution, and uses backward uniqueness for vorticity together with unique
continuation to rule out that limit.

Its implication for the fixed maximal history is exact.  Suppose
\[
 \limsup_{t\uparrow T_*}\norm{u(t)}_3<\infty.
\]
The finite limsup bounds \(\norm{u(t)}_3\) on some terminal interval.  On the
earlier compact interval the classical solution is continuous into \(H^3\)
and hence bounded in \(L^3\).  Thus
\[
 u\in L^\infty(0,T_*;L^3).
\]
By \eqref{eq:weak-strong-identification}, the suitable Leray--Hopf solution
\(U\) has the same membership and agrees with \(u\) before \(T_*\).  The
endpoint theorem makes \(U\) regular through \(T_*\), and weak--strong
uniqueness gives a classical continuation of \(u\).  Maximality rules this
out.  We have proved the necessary condition
\begin{equation}
 T_*<\infty
 \quad\Longrightarrow\quad
 \limsup_{t\uparrow T_*}\norm{u(t)}_{L^3}=\infty.
 \label{eq:ESS-blowup}
\end{equation}

The conclusion says that the critical norm crosses arbitrarily high levels
along times approaching \(T_*\).  By itself it allows a sequence of high peaks
separated by returns below a fixed level.  That distinction matters for any
argument that needs a full terminal limit rather than a selected sequence.

\subsection{Seregin's terminal divergence and Tao's quantitative rate}

Seregin closed the difference between a limsup and a limit in the smooth,
compactly supported, divergence-free whole-space class stated in his theorem.
For a solution in that class, if \(T_*<\infty\), then
\begin{equation}
 \lim_{t\uparrow T_*}\norm{u(t)}_{L^3}=\infty.
 \label{eq:seregin-L3-divergence}
\end{equation}
Equivalently, for every \(M<\infty\) there is a time \(t_M<T_*\) such that
\(\norm{u(t)}_3>M\) for every \(t\in(t_M,T_*)\)\bibref{Seregin2012}.
A finite singular branch cannot keep returning to a
fixed bounded \(L^3\) level.

Compactly supported smooth data form a proper subfamily of the class
\(\SSigma(\R^3)\) fixed in the opening.  Within that subfamily, Seregin's
theorem replaces arbitrarily late peaks by permanent terminal escape.  For
the full class, the all-data endpoint condition \eqref{eq:ESS-blowup} remains
the available conclusion.

The homogeneous Sobolev norm
\[
 \norm{u}_{\dot H^{1/2}}
 =\norm{\Lambda^{1/2}u}_2
\]
has the same scaling as \(L^3\), and the Sobolev embedding
\(\dot H^{1/2}(\R^3)\hookrightarrow L^3(\R^3)\) places it at the same
critical height.  In the same compact-data setting, Seregin proved the direct
terminal condition
\begin{equation}
 T_*<\infty
 \quad\Longrightarrow\quad
 \lim_{t\uparrow T_*}\norm{u(t)}_{\dot H^{1/2}}=\infty.
 \label{eq:seregin-critical-divergence}
\end{equation}
This is a statement about the whole terminal approach, not a
subsequence\bibref{SereginCritical}.

These theorems specify that the critical norm must diverge, but a qualitative
limit does not say how rapidly it must do so.  Tao quantified the endpoint
argument by replacing its compactness and uniqueness steps with explicit
estimates.  The change of variables
\[
 v(x,s)=\nu^{-1}u(x,s/\nu),
 \qquad q(x,s)=\nu^{-2}p(x,s/\nu),
 \qquad S_*=\nu T_*,
\]
converts the viscosity-\(\nu\) equation into the viscosity-one equation.
A further Navier--Stokes scaling normalizes \(S_*\) to one without changing
the \(L^3\) norm.  Applied after these two normalizations, Tao's estimate gives
an absolute \(c>0\) such that
\[
 \limsup_{t\uparrow T_*}
 \frac{\nu^{-1}\norm{u(t)}_3}
 {\left(\log\log\log\left(1000+
   \dfrac{T_*}{T_*-t}\right)\right)^c}
 =\infty.
\]
The time ratio and the normalized amplitude are dimensionless
\bibref{Tao2021Quantitative}.  Since
\(\dot H^{1/2}(\R^3)\hookrightarrow L^3(\R^3)\), the same lower-rate
conclusion transfers, up to the Sobolev constant, to excursions of the
critical Sobolev norm.  The estimate is slow, but it makes the shrinking-time
picture quantitative: as the remaining time collapses, critical escape cannot
stay below every power of this triple logarithm along all terminal sequences.

\Needspace{12\baselineskip}
\section{The question left by the classical theory}
\label{sec:classical-question}

The historical picture is now tight enough to state the remaining problem.
Energy allows increasing derivative intensity only by concentrating it into
less time.  Caffarelli--Kohn--Nirenberg retain scale-invariant activity near
any singular point.  ESS and Tao force arbitrarily late critical excursions;
Seregin forces permanent terminal escape in his compact-data class.  Each
result describes what a finite singular branch would have to do.  None gives
the upper estimate that would stop it.  The missing statement must be generated
by \(u_0\), \(\nu\), and the coupled equation, and its constant must remain
finite all the way to \(T_*\).

Continuation now has two distinct burdens.  Following the solution forward,
one needs an estimate strong enough to prevent its \(H^3\) restart norm from
escaping before \(T_*\).  At the endpoint, any proposed limits for the velocity
and its required derivatives must agree wherever they overlap and must still
satisfy the original momentum and normalized pressure laws.  We call the
forward estimate Gold and the endpoint test Silver.  Both belong to the same
regularity claim.

\subsection{Gold: direct continuation control}
\label{sec:gold-standard}

Fix one exponent \(s_0>5/2\).  A decisive estimate would have the form
\begin{equation}
 \sup_{0\le t<\min\{T,T_*\}}\norm{u(t)}_{H^{s_0}}
 \le C_{s_0}(u_0,\nu,T)<\infty
 \qquad\text{for every finite }T>0.
 \label{eq:gold-estimate}
\end{equation}
The constant may depend on the finite time and on the chosen Sobolev order.
It must be generated by \(u_0\), \(\nu\), and the equation; the finiteness of
the norm on the left cannot reappear as an assumption hidden inside the
constant.  Estimate~\eqref{eq:gold-estimate} is the \emph{Gold Standard}.

The threshold \(s_0>5/2\) makes the estimate continuation-grade because
\begin{equation*}
 H^{s_0}(\R^3)\hookrightarrow W^{1,\infty}(\R^3).
\end{equation*}
For every finite \(T\), it follows that
\begin{equation*}
 \int_0^{\min\{T,T_*\}}\norm{\nabla u(t)}_\infty\,dt<\infty.
\end{equation*}
Equation~\eqref{eq:high-order-gronwall} then propagates every fixed finite
Sobolev order present in \(u_0\).  The bound may change with the order.  What
smoothness needs is one finite bound at every requested order on the same time
interval, not one numerical constant for the infinite family.

Once the spatial orders are available, pressure follows from
\eqref{eq:pressure}, and the momentum equation gives
\begin{equation*}
 u_t=\nu\Delta u-(u\cdot\nabla)u-\nabla p,
\end{equation*}
so repeated differentiation recovers every finite mixed velocity--pressure
derivative.  If \(T_*<\infty\), take \(T=T_*\) in
\eqref{eq:gold-estimate}.  The propagated \(H^3\) bound gives every sufficiently
late slice the same positive local lifespan.  One such restart crosses
\(T_*\), and uniqueness identifies it with the original solution on the
overlap.  Thus the estimate alone forces \(T_*=\infty\).

Gold has now reduced the global problem to one forward task: derive
whole-terminal continuation control from the coupled equation without assuming
the norm being controlled.  The same implication can also be read from its
other end.

\subsection{Silver: testing the endpoint against the original problem}
\label{sec:silver-standard}

Silver begins by defining admissibility without placing smoothness inside the
definition.  A candidate terminal description consists of the fixed classical
pair on \([0,T_*)\) together with proposed endpoint distributions.  Call that
description admissible when it extends the same datum, the same positive
viscosity, zero forcing, incompressibility, the momentum equation, the
normalized pressure, and finite energy through \(T_*\), while agreeing with
the classical pair before the endpoint.  The target implication says that
every admissible description is smooth.  Its contrapositive says that a
genuinely nonsmooth candidate must fail at least one of those original
requirements.  This is the Silver reading.

The definition has not yet supplied a test.  Endpoint values alone cannot say
whether the momentum and pressure laws still hold.  We must first see how the
equation acts on the moving fluid, and then distinguish the ways that action
could fail at a terminal time.

\subsection{The equation acting on the fluid}
\label{sec:membership-participation}

The equations show how those requirements act in the motion.  Let \(X(a,t)\)
be the trajectory issuing from \(a\):
\begin{equation*}
 \dot X(a,t)=u(t,X(a,t)),
 \qquad X(a,0)=a.
\end{equation*}
The derivative seen while following the fluid is
\begin{equation*}
 D_t:=\partial_t+u\cdot\nabla.
\end{equation*}
The momentum equation becomes the material acceleration law
\begin{equation*}
 D_tu=-\nabla p+\nu\Delta u.
\end{equation*}
The velocity of a moving particle changes through the pressure and viscous
response of the field that contains it.  This is not an additional condition
placed on the motion; it is what the momentum equation says the motion is.

The two terms on the right form one stress divergence.  With
\begin{equation*}
 S:=\frac12\bigl(\nabla u+(\nabla u)^{\mathsf T}\bigr),
\end{equation*}
incompressibility gives
\begin{equation*}
 \nabla\cdot(2\nu S)
 =\nu\Delta u+\nu\nabla(\nabla\cdot u)
 =\nu\Delta u.
\end{equation*}
Consequently,
\begin{equation*}
 D_tu=\nabla\cdot(-pI+2\nu S).
\end{equation*}
The same transported velocity supplies both the local viscous curvature and
the velocity products from which pressure is recovered.

Jacobi's formula shows what incompressibility preserves along this motion:
\begin{equation*}
 \frac d{dt}\det\nabla_aX(a,t)
 =\bigl(\nabla\cdot u\bigr)(t,X(a,t))
  \det\nabla_aX(a,t)=0.
\end{equation*}
\Needspace{12\baselineskip}
Material volume is preserved, while \eqref{eq:pressure} reconstructs pressure
from the velocity over the whole field.  The equation joins three
scales of one event: transport follows the moving particle, viscosity reads
local spatial curvature, and pressure enforces the global incompressibility
constraint.  A candidate endpoint cannot be tested by its velocity value
alone; it must retain this response relation.  We use
\emph{viscosity--pressure--incompressibility response}, abbreviated VPI, for
this joined action on the transported velocity.  VPI is shorthand for the
original momentum, pressure, and divergence equations acting together; it is
not an extra equation.  Removing viscosity while leaving transport, pressure,
and incompressibility in place makes the first possible loss visible.

\Needspace{14\baselineskip}
\subsection{A tangential slip sheet at the inviscid boundary}
\label{sec:euler-slip}

Euler retains incompressibility, nonlinear transport, and pressure but removes
the viscous term that reads tangential curvature.  Across the interface
\(\{y=0\}\), fix constants \(a\ne b\) and define
\begin{equation}
 u(x,y,z,t)=
 \begin{cases}
  (a,0,0),&y>0,\\
  (b,0,0),&y<0,
 \end{cases}
 \qquad p(x,y,z,t)=p_0,
 \label{eq:euler-slip-sheet}
\end{equation}
where \(p_0\) is constant.  Both one-sided velocities are finite.  Their
normal components agree, while their tangential components do not.

Let \(n=(0,1,0)\).  Since
\begin{equation*}
 [u\cdot n]=0,
\end{equation*}
the distributional divergence has no sheet contribution.  In conservative
form, the possible interface defect in the Euler momentum law is
\begin{equation*}
 [(u\otimes u+pI)n]
 =[u(u\cdot n)+pn]
 =[p]n=0.
\end{equation*}
Together with \(u_t=0\), this proves that
\eqref{eq:euler-slip-sheet} is a stationary weak Euler solution.

The same field fails the fixed-viscosity equation.  Writing
\begin{equation*}
 u_1=b+(a-b)H(y),
\end{equation*}
gives
\begin{equation*}
 \partial_yu_1=(a-b)\delta_{y=0},
 \qquad
 \partial_{yy}u_1=(a-b)\delta'_{y=0}.
\end{equation*}
The viscous term \(\nu\Delta u\) contains the nonzero distribution
\(\nu(a-b)\delta'_{y=0}e_1\).  The stationary transport term supplies no
matching distribution.  Pressure cannot absorb it either: the curl of this
viscous defect is nonzero, whereas the curl of every pressure gradient is
zero.  The sheet lies on the Euler side of the fixed-viscosity
boundary.

This infinite sheet is not a finite-energy endpoint generated by the datum in
the main theorem.  Its role is narrower: it isolates the difference between
two finite one-sided values and one compatible viscous state.  Euler admits
the two tangential traces, while fixed positive viscosity detects their
distributional curvature and rejects the sheet.  A terminal failure could
postpone the same incompatibility to a gradient, acceleration, or higher
derivative even when \(u\) itself has one limit.  The endpoint argument must
produce one trace at every required finite order; finite readings along
separate approaches do not join themselves.

\Needspace{8\baselineskip}
\subsection{What can be lost at the endpoint}
\label{sec:failure-lenses}

The slip sheet and the classical continuation criteria leave three possible
endpoint failures.  A
continuation-grade quantity can become unbounded (\emph{Blown}); finite
one-sided readings can fail to join (\emph{Jump}); or the limiting values can
cease to satisfy the pressure--viscosity response (\emph{Dead}).  These are
three obligations of one terminal state---finiteness, compatibility, and the
original equation---not three separate solutions.

\paragraph{Dead.}
A Dead state has lost lawful response.  It is not a zero velocity or a zero
derivative: the zero solution, a stagnation point, and a steady flow can all
satisfy Navier--Stokes exactly.  At the base rung the test is
\begin{equation*}
 D_tu+\nabla p-\nu\Delta u=0.
\end{equation*}
If a proposed endpoint assigns a material acceleration that does not equal the
limiting pressure--viscosity response of its surrounding field, that endpoint
is Dead at the base rung.  The same mismatch can occur higher when a temporal
row no longer answers the transport, diffusion, and pressure terms obtained by
differentiating the equation.  Excluding Dead requires more than a
terminal velocity: the endpoint rows and pressure must satisfy the complete
differentiated law and generate the same local restart.

\paragraph{Jump.}
A Jump is an incompatibility between limiting values of the velocity or of a
finite derivative across a point, interface, or time.  The tangential slip
sheet \eqref{eq:euler-slip-sheet} is the interface example.  At a terminal
time, the same issue can occur when two regular approaches assign different
limits to one spatial or temporal derivative.  Both limits may be finite.
Separate bounds do not make them equal.  Excluding Jump requires a topology in
which the evolving derivative has one continuous representative and hence one
terminal value.  Part~II will derive that topology from the equation's
parabolic separation between one time derivative and two spatial derivatives.

\paragraph{Blown.}
A Blown state has an unbounded velocity, spatial derivative, temporal
response, pressure response, or other finite-depth tower readout.  The velocity
itself need not diverge.  The energy identity can remain intact while
enstrophy or a stronger continuation norm escapes.  Excluding Blown requires a
whole-terminal finite bound for every fixed temporal row and every fixed
spatial order needed by smoothness.  The constants may depend on those finite
orders; no norm uniform over the infinite tower is required.

Restart gives these failures their practical significance.  At every regular
time \(t_0<T_*\), the value \(u(t_0)\) is admissible initial data for the local
construction in Proposition~\ref{prop:maximal-history}.  A continuation-grade
bound gives a positive lifespan depending only on that bound and \(\nu\).  A
late enough slice then generates a solution beyond the candidate endpoint.
Regularity is not a list of finite quantities.  It is enough coherent state
to let the same equation run again.

\Needspace{12\baselineskip}
\section{What a finite-time singularity demands from the response}
\label{sec:singularity-tower-preview}

Dead, Jump, and Blown say what a successful endpoint argument must exclude.
They do not yet identify the response that could carry a loss toward that
endpoint.  A finite singular time need not send the velocity itself to
infinity: the energy can remain bounded while a spatial derivative, a pressure
response, or another continuation-grade quantity escapes.  To follow that
change, we return to the acceleration law of the moving fluid.

Along one trajectory, set
\begin{equation*}
 A_j(t,a):=D_t^ju(t,X(a,t)).
\end{equation*}
Since differentiation along \(X\) is \(D_t\), the fundamental theorem of
calculus gives, for every \(j\ge0\),
\begin{equation*}
 A_j(t,a)-A_j(r,a)
 =\int_r^t A_{j+1}(s,a)\,ds.
\end{equation*}
If the carried velocity escapes before a finite time, its acceleration must
accumulate without a finite \(L^1\) bound.  If the acceleration escapes, its
material jerk must carry that change.  The same implication repeats through
\begin{equation*}
 u,\qquad D_tu,\qquad D_t^2u,\qquad D_t^3u,\quad\ldots.
\end{equation*}

This is the physical reason to look beyond velocity, but it is not yet a
whole-field estimate.  A Sobolev norm measures the entire velocity field at
one time; its growth does not select a particle whose velocity or acceleration
diverges.  Spatial redistribution can enlarge that norm while every proposed
trajectory remains the wrong witness.  The next calculation must begin from
the critical quantity that the classical theory already forces to escape,
identify the fixed-coordinate response carrying its change, and retain the
transport, viscosity, and normalized pressure terms created in passing from
the moving fluid to the whole field.
